\subsection*{MT-NB01. Quantitative approximation after the adopted Riemann hypothesis}

The qualitative Nyman--Beurling and B\'aez-Duarte criteria are
already background in \texttt{RETURN\_ALPOGE\_BRIDGES\_cp02},
CC-P01--CC-P09. Here the additional issue is the rate, including
the contribution of multiple zeros. Put
\[
 d_N^2=\inf_{a_1,\ldots,a_N\in\mathbf C}
 \frac1{2\pi}\int_{\mathbf R}
 \left|1-\zeta(\tfrac12+it)
       \sum_{n\le N}a_n n^{-1/2-it}\right|^2
                 \frac{dt}{1/4+t^2}.
 \tag{NB01.1}
\]
The adopted input $A_D$ and the classical theorems below give,
for every $\epsilon>0$,
\[
 d_N^2\ll_\epsilon
 \frac{(\log\log N)^{5/2+\epsilon}}{\sqrt{\log N}},
 \qquad
 \liminf_{N\to\infty}d_N^2\log N
       \ge C_\zeta^{(2)}
       :=\sum_{\rho\ {
m distinct}}
                    \frac{m_\rho^2}{|\rho|^2}.
 \tag{NB01.2}
\]
The upper bound is Balazard--de Roton,
\emph{Sur un crit\`ere de B\'aez-Duarte pour l'hypoth\`ese
de Riemann}, Int. J. Number Theory 6(4) (2010), 883--903,
unnumbered main Theorem in \S1; author version
\texttt{arXiv:0812.1689v1}, p.~3. Its constructive choice
appears in \S2, Propositions~1--3, author pp.~4--5:
one can take
\[
 a_n=\mu(n)n^{-\varepsilon_N},\qquad
 \varepsilon_N=25(\log\log N)^{5/2+\epsilon}
                         (\log N)^{-1/2}
 \tag{NB01.3}
\]
for all sufficiently large $N$. The implied constants are
effective in the source's convention; no numerical constant
or numerical threshold is extracted here.
The lower bound is Burnol, \emph{A lower bound in an
approximation problem involving the zeros of the Riemann
zeta function}, Adv. Math. 170 (2002), 56--70,
Theorem~1.3, author version \texttt{math/0103058v2}, p.~4.
That theorem uses continuous dilations $\lambda\le\alpha\le1$;
the discrete space in (NB01.1) is a subspace at $\lambda=1/N$,
so its distance is at least the continuous distance.
The Mellin--Plancherel identification is stated and proved in
Balazard--de Roton, Proposition~1, author pp.~4--5.

The familiar sum over zeros with their ordinary multiplicities
is, under $A_D$,
\[
 C_\zeta^{(1)}:=\sum_\rho\frac{m_\rho}{|\rho|^2}
                    =2+\gamma-\log(4\pi).
 \tag{NB01.4}
\]
An exact locator for the identity in this normalization is
Balazard--de Roton, author p.~3, immediately before (3),
and Bettin--Conrey--Farmer, author p.~2.
Thus $C_\zeta^{(2)}\ge C_\zeta^{(1)}$, with equality precisely
when every nontrivial zero is simple. The adopted RH input
does not remove this distinction.

For comparison, Bettin--Conrey--Farmer,
\emph{An optimal choice of Dirichlet polynomials for the
Nyman--Beurling criterion}, Proc. Steklov Inst. Math.
280, Suppl.~2 (2013), S30--S36, Theorem~1,
\texttt{arXiv:1211.5191v1}, p.~2, proves the sharper conclusion
\[
 d_N^2\sim\frac{2+\gamma-\log(4\pi)}{\log N}
 \tag{NB01.5}
\]
under RH and the additional hypothesis
\[
 \sum_{|\operatorname{Im}\rho|\le T}
                    |\zeta'(\rho)|^{-2}
                  \ll T^{3/2-\delta}\quad\text{for some }\delta>0.
 \tag{NB01.6}
\]
The hypothesis includes simplicity. Their explicit polynomial is
$V_N(s)=\sum_{n\le N}(1-\log n/\log N)\mu(n)n^{-s}$.
The theorem gives this polynomial's error asymptotic; Burnol's
lower bound then gives (NB01.5) for the infimum. These are
classical results with their exact remaining inputs, not new
proofs of the cited analytic estimates.

\subsection*{MT-NB02. The multiplicity-sensitive Hecke approximation constant}

Let $K$ be a number field, and let $F(s)=L_K(s,\eta)$ be a
primitive unitary Hecke $L$-function. Allow $\eta=1$, but exclude
the nontrivial pure norm characters $|\cdot|^{it}$, whose pole
would lie away from $1$. For such a pure norm character the
following statement instead applies to its vertical translate
$\zeta_K(s)$. Write $F(s)=\sum_{n\ge1}a_n n^{-s}$ and define
\[
 \Psi_F(x)=\operatorname*{Res}_{s=1}
                   \frac{x^s F(s)}s-\sum_{n\le x}a_n,
 \qquad
 B_F^\lambda=\operatorname{span}\{
                  t\longmapsto\Psi_F(\alpha/t):
                               \lambda\le\alpha\le1\}.
 \tag{NB02.1}
\]
Assume GRH for this particular $F$, supplied by $A_D$ for a
primitive Dirichlet function and by the conditional input
$M_H$ in the general Hecke case. Then $\Psi_F(1/t)$ belongs
to $L^2((0,\infty),dt)$, the distance
\[
 D_F(\lambda)=\operatorname{dist}_{L^2(0,\infty)}
                      (\mathbf1_{(0,1)},B_F^\lambda)
 \tag{NB02.2}
\]
tends to zero, and
\[
 \liminf_{\lambda\downarrow0}D_F(\lambda)^2\log(1/\lambda)
    \ge C_F^{(2)}:=\sum_{\rho\ {
m distinct}}
                          \frac{m_\rho^2}{|\rho|^2}<\infty.
 \tag{NB02.3}
\]
The zeros in this formula are the nontrivial zeros of $F$,
all on the critical line under the stated input.
The precise classical theorem being specialized is
Anne de Roton, \emph{Une approche hilbertienne de
l'hypoth\`ese de Riemann g\'en\'eralis\'ee}, Bull. Soc.
Math. France 134(3) (2006), 417--445: Definition~1 p.~418,
Definition~2 and Theorem~1 p.~419, Theorem~2 p.~420.
The lower-limit argument is \S7, pp.~441--444, particularly
Propositions~7.1--7.2 and \S7.2. The original p.~420 image
and that complete final proof were read. Earlier technical
operator estimates remain cited prerequisites.

\begin{proof}[Complete proof of the stated Hecke specialization]
We check the source's Selberg-class hypotheses, rather than
assuming that an arbitrary meromorphic Artin quotient satisfies
them. Grouping the ideal Dirichlet series by absolute norm gives
the displayed coefficients. If $n_K=[K:\mathbf Q]$, their
absolute values are bounded by $d_{n_K}(n)$: at each rational
prime the local ideal-counting coefficients are bounded by those
of $(1-p^{-s})^{-n_K}$. Consequently $a_n\ll_\epsilon n^\epsilon$.
The unitary Euler factors have roots of modulus at most one.
For a rational prime $p$, the logarithmic coefficient at $p^r$
is a sum over prime ideals of residue degree $f\mid r$,
with each summand bounded by $f/r$. Since the sum of the
residue degrees is at most $n_K$, its absolute value is at
most $n_K/r$. All other logarithmic coefficients vanish.
This verifies the source's Euler-product bound with exponent
zero, strictly smaller than $1/2$.

The classical Hecke continuation, finite-order property and
functional equation supply the remaining axioms; see the exact
Ramakrishnan--Valenza Theorem~7-19 and local-factor locators
in MT-M28. The real gamma shifts have nonnegative real part,
and the complex shifts have real part $|m|/2\ge0$.
There is no pole for a non-pure-norm character. The allowed
trivial character gives $\zeta_K$, with its single pole at $1$.
Thus the permitted functions belong to the precise Selberg
class of de Roton's Definition~1. Her Theorem~1, with the
specified GRH, gives the $L^2$ assertion and convergence of
the distance defined in (NB02.1)--(NB02.2). Her Theorem~2
then gives the lower bound; squaring a nonnegative lower-limit
inequality yields (NB02.3).

For clarity, convergence of the multiplicity-square sum also
follows directly from the fixed-function local zero bound in
NT-Z10. In every unit height interval the sum of multiplicities
is $O_F(\log(2+|T|))$. The sum of their squares is at most
the square of that sum. The contribution of intervals of
height at least two is therefore bounded by a constant times
$\sum_{j\ge2}j^{-2}\log^2j$, which converges. The finitely many
bounded-height zeros cause no problem, since a critical-line
zero has $|\rho|\ge1/2$.
\end{proof}

\subsection*{MT-NB03. Integral Galois classes force an approximation penalty}

For a function in MT-NB02 set
\[
 C_F^{(1)}=\sum_\rho\frac{m_\rho}{|\rho|^2},\qquad
 \Delta_F=\sum_\rho\frac{m_\rho(m_\rho-1)}{|\rho|^2}.
 \tag{NB03.1}
\]
Then $C_F^{(2)}=C_F^{(1)}+\Delta_F$, and
$\Delta_F=0$ precisely when all its nontrivial zeros are simple.
For a tame finite Galois extension $N/K$ put
$U_{N/K}=[\mathcal O_N]-[\mathcal O_K[G]]$
in the locally free class group of $\mathbf Z[G]$, with the
restriction-of-scalars convention of MT-ART11. If
$U_{N/K}\ne0$, then, conditionally on GRH for $\zeta_N$,
\[
 \begin{split}
 \liminf_{\lambda\downarrow0}D_{\zeta_N}(\lambda)^2
                      \log(1/\lambda)&\ge C_{\zeta_N}^{(2)},\\
 C_{\zeta_N}^{(2)}&\ge16,\qquad
 C_{\zeta_N}^{(2)}\ge C_{\zeta_N}^{(1)}+8.
 \end{split}                                                   \tag{NB03.2}
\]
The same constants apply to the explicit tame quaternionic
field of MT-ART10, directly from its proved central order
$r\ge2$; this use does not identify nonfreeness with a nonzero
stable class. Independently of $U_{N/K}$, if $N/K$ is nonabelian
Galois, then $\Delta_{\zeta_N}>0$ under that same GRH input.
Thus a multiplicity-one constant is strictly too small in
all these nonabelian examples.

\begin{proof}
All the sums converge by the argument in MT-NB02. The identity
in (NB03.1) follows term by term from $m^2=m+m(m-1)$.
Every summand of $\Delta_F$ is nonnegative, and it vanishes
exactly when its positive integer multiplicity equals one.

MT-ART11 proves unconditionally, without Artin holomorphy,
that $U_{N/K}\ne0$ forces an even positive central order
$r=\operatorname{ord}_{s=1/2}\zeta_N(s)\ge2$. Its proof uses
Taylor's root-number class theorem and the cyclic-restriction
argument recorded there. The single central zero contributes
$4r^2\ge16$ to $C_{\zeta_N}^{(2)}$ and
$4r(r-1)\ge8$ to $\Delta_{\zeta_N}$. This proves (NB03.2)
by MT-NB02. The arithmetic central-zero assertion is
unconditional; GRH enters to obtain the approximation-space
conclusion and place every nontrivial zero in that formula.

For any nonabelian $N/K$, the classical
Hu--Kaneko--Martin--Schildkraut Theorem~1.1, already located
and used in MT-M23, gives infinitely many multiple nontrivial
zeros of $\zeta_N$. Under the present GRH each contributes a
strictly positive term to $\Delta_{\zeta_N}$; even one would
suffice. Hence the claimed strict inequality follows. No
uniform numerical lower bound is asserted in this last case.
\end{proof}

\subsection*{MT-NB04. A precise derivative exponent sufficient for the optimal rate}

Suppose in addition to $A_D$ that, for some $c>0$ and
$0\le\beta<1/4$, every nontrivial zero satisfies
\[
             |\zeta'(\rho)|\ge c(1+|\operatorname{Im}\rho|)^{-\beta}.
 \tag{NB04.1}
\]
Then all the zeros are simple, (NB01.6) holds for some
$\delta>0$, and the optimal asymptotic (NB01.5) follows.

\begin{proof}
The displayed positive lower bound implies simplicity. The
classical zero-count bound $N(T)=O(T\log(2+T))$, with the
same multiplicity convention as NT-Z10, gives for $T\ge2$
\[
 \sum_{|\operatorname{Im}\rho|\le T}|\zeta'(\rho)|^{-2}
                 \ll_{c,\beta}T^{1+2\beta}\log T.
\]
Take $\delta=1/4-\beta>0$. Since $\log T\ll_\delta T^\delta$,
the right side is $O(T^{5/4+\beta})=O(T^{3/2-\delta})$.
Bettin--Conrey--Farmer Theorem~1 gives the error asymptotic
for $V_N$, and Burnol gives the matching lower bound for
the infimum, as in MT-NB01. This proves the conclusion.
\end{proof}

This sufficient exponent condition is stronger than the
unspecified polynomial reciprocal-derivative bound forced by
the point-jet basis in MT-M11. It is not a consequence here
of $A_D$, of the repaired adapted-block basis, or of simple
zero locations alone. It makes the additional quantitative
condition needed by that particular route explicit.

\paragraph{Geometric and run placement of MT-NB01--04.}
The Hecke functions are the rank-one adelic objects of Tate's
thesis; for $\zeta_N$ the object is the pushforward of the unit
Artin motive along $\operatorname{Spec}N\to\operatorname{Spec}K$.
The functions $\Psi_F$ turn its coefficient-counting distribution
into a dilation subspace. The zero divisor, including its
orders, obstructs approximation by that subspace. In de Roton's
proof each order-$m$ zero supplies a full $m$-jet Gram block;
the inverse block's first diagonal entry is $m^2$ (\S7.2,
p.~444). This is the analytic reason that discarding higher
jets loses the constants in MT-NB03, consistent with MT-M20
and MT-M25.

For the archive programme, this is an approximation-theoretic
addition to the trace/Fredholm interface of cote~1, scans
79--80 and 111--112, inherited through
\texttt{OURS\_INTAKE\_RESEARCH\_SA23\_20261120.md} \S4.1,
\texttt{geo-H-lidskii}/\texttt{geo-H-genus}; the sheets themselves
are not cited as Nyman--Beurling theorems. MT-NB03 also uses
the integral $K_0$ comparison of MT-ART11 and
AG-LEFSCH-LATTICE/DEFECT, placed at cote~16 scans~14--18,
no.~7, and SGA6 Expos\'e~V \S2 and \S7 (7.1.1), with the
same inherited-locator qualification as MT-ART11. The precise
run extension is the quantitative completion of the qualitative
CC-P01--CC-P09 bridge, with the negative derivative moment
(NB01.6) retained as the remaining optimal-rate condition.
